\chapter{Análisis de Requisitos}
\label{chap:requisitos}

En este capítulo se establecen los requisitos que la solución de software propuesta debe satisfacer. Estos se dividen en dos categorías principales: los requisitos funcionales, que describen las funcionalidades y capacidades que el sistema debe ofrecer (el <<qué>>), y los requisitos no funcionales, que definen los atributos de calidad, restricciones y características operativas del sistema (el <<cómo de bien>>).

\section{Requisitos Funcionales (RF)}
\label{sec:rf}

Los requisitos funcionales del sistema se han agrupado en tres categorías principales que se corresponden con las fases lógicas del flujo de trabajo: (1) Ingesta y Procesamiento Estructural de los Documentos; (2) Representación y Enriquecimiento del Conocimiento; y (3) Interacción, Búsqueda y Generación de Respuestas. A continuación, se detalla cada grupo en tablas separadas.
\newpage
\subsubsection*{Grupo 1: Ingesta y Procesamiento Estructural de Documentos}
Estos requisitos describen las funcionalidades necesarias para cargar, analizar y segmentar los documentos de entrada, transformando el contenido no estructurado en una representación jerárquica inicial.

\begin{table}[H]
\centering
\renewcommand{\arraystretch}{1.4} % Aumenta el espacio vertical para mejorar la legibilidad
\caption{Requisitos Funcionales: Ingesta y Procesamiento Estructural.}
\label{tab:rf_ingesta}
\begin{tabularx}{\textwidth}{|l|l|X|}
\hline
\textbf{ID} & \textbf{Nombre} & \textbf{Descripción} \\
\hline
\textbf{RF-01} & Carga de Documentos & El sistema deberá permitir la carga simultánea de múltiples documentos en formato PDF desde una ruta de directorio especificada. \\
\hline
\textbf{RF-02} & Detección Adaptativa de TOC & El sistema deberá implementar una estrategia dual para detectar la Tabla de Contenidos (TOC): prioritariamente mediante el análisis de hipervínculos internos y, como alternativa, mediante el reconocimiento de patrones textuales con expresiones regulares. \\
\hline
\textbf{RF-03} & Inferencia de Jerarquía & El sistema deberá construir una estructura jerárquica (padre-hijo) de las secciones del documento, basándose en el análisis de prefijos numéricos (ej: 1.2.1) y la indentación visual (posición X) de los títulos en la TOC. \\
\hline
\textbf{RF-04} & Chunking Jerárquico & El sistema deberá segmentar (chunking) los documentos de forma estructural, creando chunks cuyo contenido corresponde al texto delimitado entre el título de una sección y el de la siguiente, localizando dichos títulos con precisión mediante similitud difusa (fuzzy matching). \\
\hline
\end{tabularx}
\end{table}

\newpage
\subsubsection*{Grupo 2: Representación y Enriquecimiento del Conocimiento}
Una vez procesada la estructura inicial, los siguientes requisitos se centran en cómo esta información es representada, almacenada y enriquecida semánticamente para facilitar una recuperación de información más avanzada.

\begin{table}[H]
\centering
\renewcommand{\arraystretch}{1.4}
\caption{Requisitos Funcionales: Representación y Enriquecimiento.}
\label{tab:rf_representacion}
\begin{tabularx}{\textwidth}{|l|l|X|}
\hline
\textbf{ID} & \textbf{Nombre} & \textbf{Descripción} \\
\hline
\textbf{RF-05} & Construcción de Grafo Estructural & Cada documento deberá ser representado como un grafo dirigido (\texttt{NetworkX}), donde los nodos son los chunks y las aristas representan relaciones estructurales (jerárquicas y secuenciales). \\
\hline
\textbf{RF-06} & Enriquecimiento Semántico del Grafo & El grafo estructural deberá ser enriquecido mediante la adición de aristas semánticas entre nodos no conectados estructuralmente que superen un umbral predefinido de similitud de coseno entre sus embeddings. \\
\hline
\textbf{RF-07} & Generación de Embeddings Duales & Para cada chunk, el sistema deberá generar y almacenar dos representaciones vectoriales (embeddings) distintas: una para su título y otra para su contenido textual. \\
\hline
\end{tabularx}
\end{table}

\newpage
\subsubsection*{Grupo 3: Interacción, Búsqueda y Generación (GraphRAG)}
Este grupo final de requisitos define el núcleo interactivo del sistema, detallando cómo se procesa la consulta del usuario, se recupera la información relevante del grafo de conocimiento y se genera la respuesta final.

\begin{table}[H]
\centering
\renewcommand{\arraystretch}{1.4}
\caption{Requisitos Funcionales: Interacción, Búsqueda y Generación.}
\label{tab:rf_interaccion}
\begin{tabularx}{\textwidth}{|l|l|X|}
\hline
\textbf{ID} & \textbf{Nombre} & \textbf{Descripción} \\
\hline
\textbf{RF-08} & Selección de Estrategia por Agente & Un agente basado en LLM deberá analizar la consulta del usuario para determinar la estrategia de respuesta óptima: Acceso Parcial (RAG) para preguntas específicas, o Acceso Completo para preguntas generales (ej: resúmenes). \\
\hline
\textbf{RF-09} & Optimización de Consulta por Agente & Un agente LLM deberá refinar la consulta del usuario antes de la búsqueda, simplificándola a sus términos clave para mejorar la precisión del sistema de recuperación de información. \\
\hline
\textbf{RF-10} & Búsqueda Híbrida y Re-ranking & El sistema realizará una búsqueda híbrida combinando similitud semántica (embeddings) y búsqueda lexicográfica (BM25) sobre sub-chunks. Los resultados serán reordenados mediante un modelo Cross-Encoder para maximizar la relevancia. \\
\hline
\textbf{RF-11} & Selección de Contexto Basada en Grafo & El contexto final se seleccionará analizando un subgrafo de consulta, aplicando una puntuación final a cada nodo que combina su relevancia directa, la relevancia propagada de sus familiares (padre, hijos, hermanos) y métricas de centralidad del grafo. \\
\hline
\textbf{RF-12} & Modos de Respuesta Múltiples & El sistema deberá ser capaz de generar respuestas en dos modos: individual (una respuesta por documento) y consolidado (una única respuesta sintetizada a partir de todos los documentos). \\
\hline
\end{tabularx}
\end{table}
\newpage

\section{Requisitos No Funcionales (RNF)}
\label{sec:rnf}

Los requisitos no funcionales definen los atributos de calidad del sistema, tales como su rendimiento, modularidad, robustez y mantenibilidad. Estos requisitos son cruciales para asegurar que el sistema no solo funcione, sino que lo haga de manera eficiente y fiable. A continuación, se detallan agrupados por su naturaleza.

\subsubsection*{Grupo 1: Atributos de Arquitectura y Calidad de Software}
Este grupo define las características intrínsecas del diseño y la construcción del software que garantizan su mantenibilidad, flexibilidad y adaptabilidad a futuro.

\begin{table}[H]
\centering
\renewcommand{\arraystretch}{1.4}
\caption{Requisitos No Funcionales: Arquitectura y Calidad.}
\label{tab:rnf_arquitectura}
\begin{tabularx}{\textwidth}{|l|l|X|}
\hline
\textbf{ID} & \textbf{Nombre} & \textbf{Descripción} \\
\hline
\textbf{RNF-01} & Modularidad & La arquitectura del sistema deberá estar desacoplada, con una clara separación de responsabilidades entre los componentes de extracción (\texttt{chunk\_extractor.py}), decisión (\texttt{document\_strategy\_agent.py}) y el núcleo RAG (\texttt{main.py}), facilitando su mantenimiento y escalabilidad. \\
\hline
\textbf{RNF-02} & Configurabilidad & El sistema deberá permitir la configuración del modelo de lenguaje (LLM) a utilizar a través de parámetros, sin necesidad de modificar el código fuente principal. \\
\hline
\textbf{RNF-03} & Portabilidad (Ejecución Local) & El sistema deberá ser capaz de operar utilizando modelos de lenguaje alojados localmente mediante una interfaz como Ollama, garantizando su funcionamiento sin dependencia de servicios cloud de pago. \\
\hline
\end{tabularx}
\end{table}

\newpage
\subsubsection*{Grupo 2: Atributos de Operación y Rendimiento}
Estos requisitos se centran en el comportamiento observable del sistema durante su ejecución, incluyendo su eficiencia, fiabilidad y la facilidad para diagnosticar su funcionamiento.

\begin{table}[H]
\centering
\renewcommand{\arraystretch}{1.4}
\caption{Requisitos No Funcionales: Operación y Rendimiento.}
\label{tab:rnf_operacion}
\begin{tabularx}{\textwidth}{|l|l|X|}
\hline
\textbf{ID} & \textbf{Nombre} & \textbf{Descripción} \\
\hline
\textbf{RNF-04} & Trazabilidad y Depuración & El sistema deberá generar artefactos para la depuración, incluyendo visualizaciones del grafo de consulta (\texttt{.png}) y archivos de texto (\texttt{.txt}) que registren el contexto exacto enviado al LLM en cada consulta. \\
\hline
\textbf{RNF-05} & Robustez & El sistema deberá gestionar de forma controlada los errores durante la lectura de documentos (ej: PDFs corruptos), notificando al usuario y continuando la operación con los archivos válidos sin interrumpir la ejecución. \\
\hline
\textbf{RNF-06} & Gestión de Recursos & La versión del sistema que opera con modelos locales deberá ser funcional en equipos de escritorio con especificaciones de hardware razonables (ej: 16 GB de RAM y GPU con 8 GB de VRAM). \\
\hline
\end{tabularx}
\end{table}
